\subsection{Unitary front end, D4 guard shear, and the S3 optimizer frontier}
\label{sec:bt1419-bt1421-frontier}

The BT1417 primitive list can be promoted from an incidence checklist to a
symbolic optical-unitary certificate.  The edge-channel coupler is
\[
H_2=\frac1{\sqrt2}\begin{pmatrix}1&1\\1&-1\end{pmatrix},
\]
the oriented latch is \(\operatorname{diag}(1,\pm1)\), the residue demux is the
four-mode Fourier transform \(F_4[j,k]=i^{jk}/2\), and a six-channel
Cs\'asz\'ar/Szilassi star analyzer may be represented by
\(F_6[j,k]=e^{2\pi ijk/6}/\sqrt6\).  Each block is unitary.  With a generic
beamsplitter mesh bound, the active stack has conservative depth
\[
1+1+6+15=23.
\]
This is a finite symbolic bound, not a waveguide mask or loss model.

The D4 quartic guard band also has an internal algebra.  The 24 guard apertures
are
\[
2\text{ atoms}\times4\text{ branches}\times3\text{ qutrit phases},
\]
and the D4 orientation lift gives \(24\cdot8=192\) tomotope tokens.  On one
atom the finite state space is \(\mathbb Z_4^{\rm branch}\times\mathbb Z_3^{\rm phase}\).
The Clifford frame supplies the uniform qutrit phase shift, whereas the guard
injection is the branch-controlled shear
\[
J:(b,p)\longmapsto (b,p+b\bmod3).
\]
The verifier checks \(J^3=1\), \(J\ne1\), and \(J\) is not a product of a D4
branch permutation with a uniform qutrit phase shift.  Thus the non-Clifford
resource is visible as a finite phase-frame shear.

Finally, the S3 synchronization problem is now stated as an optimizer frontier
under the complete physical front end.  The incumbent remains
\[
540=210_{\rm id}+330_{\rm corr}.
\]
The radius-three certificate excludes
\[
\sum_{r=1}^3\binom{39}{r}(6^r-1)=1,991,015
\]
root-fixed candidate relabelings, so any better global witness must lie at
radius at least four.  The front-end split records
\[
210=21\cdot10=42\cdot5,
\qquad
330=168+162.
\]
The remaining proof obligation is deliberately explicit: either certify
\(\sum_e y_e\le210\) for the 540 S3 skew-line constraints, or exhibit a gauge
with at least 211 identity residual edges.
